\subsubsection{Bounded ramification and linear cohomology do not force generation}

\begin{proposition}[Refuted: a bounded-ramification generation inference]
\label{D:bounded-ramification-generation-obstruction}
Let $k$ be an algebraically closed field of characteristic $2$.
There is a quotient $G$ of $\pi_1(\mathbf G_{m,k})$ with all
of the following properties.
\begin{enumerate}
\item Every finite quotient cover has nonlogarithmic
Abbes--Saito breaks at $0$ and $\infty$ at most $4/3$.
\item $G$ is small: it has finitely many continuous
homomorphisms to each fixed finite group.
\item The maximal pronilpotent quotient of every open
subgroup of $G$ is finite.
\item For every finite field $\Lambda$ and every
finite-dimensional continuous $\Lambda[G]$-module $M$,
\[
                   \dim_\Lambda H^1(G,M)\le2\dim_\Lambda M.
                                                        \tag{D.BR.1}
\]
For an open subgroup $H$, the corresponding bound is
$2[G:H]\dim_\Lambda M$.
\item $G$ is not topologically finitely generated.
\end{enumerate}
In particular, bounded ramification, smallness, and
all-characteristic rank-linear degree-one cohomology
do not together imply finite generation.
This is an obstruction to that inference; it is not
a counterexample to the original-divisor tame assertions
D.1 or D.2(ii).
\end{proposition}

\begin{proof}
For $i\ge1$, set
\[
 q_i=2^{i+1},\qquad S_i=\operatorname{PSL}_2(\mathbf F_{q_i}),
 \qquad N_i=|S_i|^i+1,\qquad
                   G=\prod_{i\ge1} S_i^{N_i}.             \tag{D.BR.2}
\]
The groups $S_i$ are pairwise nonisomorphic nonabelian
simple groups.  Proposition~\ref{D:psl2-bounded-conductor-covers}
proves, without a classification input, that each is
generated by three elements and supplies arbitrarily many
pairwise inequivalent $S_i$-covers of $\mathbf G_m$.
Their branch points remain $0,\infty$; the wild branch
has classical upper jump $1/(q_i-1)$, and the other branch
is tame.  The same proposition proves linear disjointness
of every finite collection of these covers, including
distinct scalar twists with the same simple group.
Choose $N_i$ twists for each $i$.  The corresponding
map from $\pi_1(\mathbf G_m)$ surjects onto every finite
subproduct in (D.BR.2).  Its image is compact and therefore
closed, so it is the whole product.

At these valuations the residue field is perfect.
The comparison with the classical filtration gives
nonlogarithmic break
\[
                       1+\frac1{q_i-1}\le\frac43
\]
at the wild branch and at most $1$ at the tame branch
\citep[Proposition~3.15(4), preprint p.~8]{AbbesSaito2002D}.
Every finite quotient of $G$ factors through a finite
subproduct.  The ramification subgroup above a bound
acts trivially in a compositum when it acts trivially
in each constituent, and passage to quotients preserves
that property.  Thus the same $4/3$ bound applies to
every finite quotient cover, not only the displayed
simple factors.

Fix a finite group $F$.  A nontrivial homomorphism
$S_i\to F$ is injective.  Therefore it is impossible
when $|S_i|>|F|$.  Only finitely many values of $i$ can
contribute, and each has only the finite multiplicity
$N_i$.  Every continuous homomorphism $G\to F$ hence
factors through one fixed finite subproduct.  This proves
smallness.

If $H$ is open in $G$, it contains a full tail
\[
                     K_m=\prod_{i>m}S_i^{N_i}.
\]
Writing $P_m=\prod_{i\le m}S_i^{N_i}$, we have
$H=H_0\times K_m$ for a subgroup $H_0\subset P_m$:
multiply any chosen lift of an element of $H_0$ by
an element of the contained tail to remove its tail
coordinate.  Every finite nilpotent image of $K_m$
is trivial.  Indeed, a homomorphism to a finite group
uses only finitely many coordinates, and a nonabelian
simple factor has no nontrivial nilpotent quotient.
The maximal pronilpotent quotient of $H$ consequently
factors through the finite group $H_0$.

We give the cohomology argument in all coefficient
characteristics, including characteristics dividing
$|S_i|$.  For a finite product $P=\prod_{\nu=1}^t S_\nu$,
every absolutely simple module over an algebraically
closed coefficient field is an outer tensor product
$V=\bigotimes_\nu V_\nu$ of simple modules for the
factors.  This follows by passing from the finite
group algebras to their semisimple quotients: the
radical of each tensor factor acts trivially on a
simple module, and tensor products of matrix algebras
have the stated simple modules.  Thus no Maschke
assumption is being used.

If no factor acts nontrivially, $V$ is trivial and
$H^1(P,V)=0$, since $P$ is perfect.  If exactly one
factor acts nontrivially, inflation and the perfectness
of the remaining factors identify $H^1(P,V)$ with
$H^1(S_\nu,V_\nu)$.  A cocycle is determined by its
values on three generators, so its space has dimension
at most $3\dim V_\nu$.  Coboundaries have dimension
$\dim V_\nu$, because $V_\nu^{S_\nu}=0$.  Hence
\[
                       \dim H^1(P,V)\le2\dim V.           \tag{D.BR.3}
\]
If at least two factors act nontrivially, write
$P=A\times B$ with one active factor in $A$ and at
least one in $B$, and $V=V_A\otimes V_B$.
The degree-one Hochschild--Serre sequence injects
$H^1(P,V)$ into $H^1(A,V)^B$, since $V^A=0$.
But
\[
 H^1(A,V)^B
       =\bigl(H^1(A,V_A)\otimes V_B\bigr)^B=0.
\]
Thus (D.BR.3) holds in every case.

Extension of the coefficient field commutes with
finite-group cohomology.  A composition series after
algebraic closure and the degree-one long exact
sequence now give (D.BR.3) for every finite-dimensional
module over every finite field.

For a continuous $G$-module $M$, choose $m$ so that
the full tail $K_m$ acts trivially.  Its first
cohomology with these coefficients is
$\operatorname{Hom}_{\rm cont}(K_m,M)=0$, since every
finite tail image is perfect.  Inflation therefore
identifies $H^1(G,M)$ with $H^1(P_m,M)$ and proves
(D.BR.1).  Shapiro applied to
$\operatorname{Ind}_H^G M$ gives the stated open-subgroup
bound.

Finally, if $G$ had $d$ topological generators, choose
$i\ge d$.  Their images would generate $S_i^{N_i}$.
There are at most $|S_i|^d\le|S_i|^i$ possible ordered
$d$-tuples of coordinate entries, but
$N_i=|S_i|^i+1$.  Two coordinates therefore receive
identical tuples.  Every word in the proposed
generators has equal entries at these coordinates,
so the generated subgroup is contained in a proper
diagonal subgroup.  This contradiction proves
non-generation.
\end{proof}

\begin{remark}[Scope and finite certification]
The covers in this example are wild at one original
boundary point.  Accordingly, $G$ is not asserted to
be a quotient of the original-divisor tame group
for the pair $(\mathbf P^1,\{0,\infty\})$.
Its role is to test the proposed abstract deduction
from conductor bounds and cohomology to generation.

The exact finite program
\nolinkurl{parts/D/code/check_psl2_bounded_ramification_v2.py}
checks the fields $q=4,8,16$, the three matrix generators,
the two branch orbits and stabilizers, the lower inertia
filtration, the invariant identities, and the numerical
product inequalities.  Its current output is
\nolinkurl{parts/D/output/psl2_bounded_ramification_v2.json}.
This restored version is separately identified from the
historical program, which had 306,144 assertions and
was independently rerun before the environment replacement.
The general proof is the argument above and the
preceding geometric proposition; the finite computation
does not replace it.
\end{remark}
